% IV. ALTERNATIVE EXPLANATIONS -- 500 words, Table II. NEW FILE.
%
% A full-length draft already exists in this repository at
% manuscript/drafts/draft-section-IV.md (moved from the study repository's
% tubitak/docs/ on 2026-09-13; 1,109 words, four refutations plus the
% design-based argument). It is drafted at arXiv length by the structural decision
% of 2026-08-26; condensing it to the letter's 500 words is a separate task and has
% not been done. DO NOT paste the long version in here.
%
% TABLE II IS THREE ROWS, not five, and not the four of the first revision:
%   1. Cold-discriminator.
%   2. Matcher independence -- re-attributed across two registrations. The "48 of
%      49" figure is STRUCK and must not be quoted: it is not reproducible from the
%      artifact.
%   3. Blur vs restraint -- REPLACED by a stronger test. The original was a
%      single-seed negative control; the replacement is the six-seed informative-
%      mask edge-ratio measurement. C2 reproduces the real image's edge density to
%      within 1.5% where the input asserts structure, so the suppression is
%      SELECTIVE (restraint), not blur. Verdict confirmed across all six seeds.
%
% REMOVED ROWS, and they do not come back:
%   - Mediation row: STRUCK (corrections-log entry 20; paper-roadmap already rules
%     B3 part 2 does not appear at all).
%   - Georeferencing row: REMOVED from the table and ANSWERED IN PROSE instead. The
%     answer is design-based rather than measured -- stronger than the 86% figure it
%     replaces, because it depends on the design rather than on one number. The
%     registered positional contrasts contain no pretrained comparisons, so the
%     objection does not reach them.
%
% The freed word budget is NOT reclaimed by this section; the allocation stands as
% revised.
%
% Mediation footnote: write it carefully. The registered mediation test was run and
% its reported conditional is void as stated -- say that, do not imply the test
% supports anything.
%
% Mechanistic note, ~40 words, because it explains why the result is
% matcher-independent.

% Migrated from manuscript/drafts/draft-section-IV.md (Draft 1, 2026-08-26, arXiv length,
% 1,109 words) on 2026-09-13. Prose only, transcribed without compression or correction;
% the condensation to the letter's 500 words is a separate task and has NOT been done.
% Markdown bold is carried as \textbf. Table II is not typeset; the reference is left as a
% placeholder. Values that are the Table I edge-ratio column on the input-silent mask are
% wrapped in \tabone{...}.

Four explanations compete with the account above. Three are tested and refuted; the
fourth is answered by the design itself and needs no test.
Table~\todo{II: three rows, not yet typeset} summarises; the prose gives each row the
evidence that supports it and the caveat that qualifies it.

\subsection{It is blur, not restraint}

\textbf{The objection.} The L1-only arm produces smoother output, and feature matchers
prefer smoother imagery. If so, its advantage is a mechanical consequence of blur rather
than of any learned decision about what to render.

\textbf{The test.} The invention measurement is defined on input-silent pixels---those
where the conditioning input asserts no structure. Computing the \emph{same} ratio on the
complementary mask, where the input does assert structure, separates the two explanations
exactly, because \textbf{blur suppresses edges uniformly while restraint suppresses them
conditionally.} A smoothing explanation predicts the L1-only arm below unity on both
masks. A restraint explanation predicts it near unity where the input is informative and
far below where it is not. The thresholds separating ``near unity'' from ``suppressed''
were fixed in advance at 0.80 and 0.50, reusing the bands already registered for the
original measurement rather than choosing new ones.

\textbf{The result.} Across six seeds the L1-only arm reproduces the real image's edge
density to within 1.5\% on the informative mask---0.986, range 0.980 to 0.989---while
sitting at \tabone{0.277} on the input-silent mask. \textbf{A factor of 3.6 between the two
masks, in the same arm, on the same chips, with the same operator, in every seed.} Both
conditions hold six times out of six. The other three arms sit between 1.03 and 1.05 on
the informative mask, slightly above reality, consistent with the \tabone{1.07 to 1.16}
they show where the input says nothing.

\textbf{The reading.} Every arm except the L1-only one adds edges everywhere. The L1-only
arm adds them only where the input warrants it. \textbf{That is conditional suppression,
which blur cannot produce}, and the objection is refuted by a positive test rather than by
a negative control.

\textbf{What this row is not.} It does not establish that the L1-only arm is optimal, nor
that restraint is the only mechanism in play---Section III's honest-limit paragraph
stands. It establishes that the smoothing explanation, specifically, does not fit the
data.

\subsection{Cold-started discriminator damage}

\textbf{The objection.} The published generator is deposited without its discriminator,
so every adversarial arm here begins from a randomly initialised one. Adversarial arms
might therefore lose because of a damaged start rather than because of the adversarial
objective.

\textbf{The test.} A checkpoint sweep at epochs 1, 2, 5, 10 and 20, comparing each arm to
the pretrained generator and to its own counterpart.

\textbf{The result.} At epoch 1 the adversarial+L1 arm is \textbf{already better than
pretrained}, by $-0.399 \pm 0.064$~px at 6.3 standard errors. \textbf{That is the wrong
sign for startup damage}: an arm crippled by a cold discriminator would begin worse than
the checkpoint it started from, not better. The adversarial deficit relative to the
non-adversarial arm is present from epoch 1 at $+0.546 \pm 0.048$ and reaches $+0.700$ by
epoch 20.

\textbf{The caveat, and it is ours rather than a reviewer's.} The path from epoch 1 to
epoch 20 is not monotone: it dips to $+0.384$ at epoch 5 before growing. \textbf{The
registered bands for this sweep did not anticipate that shape, and the conclusion is read
off the curve rather than matched to a band.} We state it that way because presenting it
as a clean band hit would misrepresent how it was obtained. The refutation does not depend
on the shape---it depends on the epoch-1 magnitude, which is large and correctly signed.

\subsection{Optimising the evaluation metric}

\textbf{The objection.} The result might be an artefact of the KLT matcher used to
produce it.

\textbf{The evidence, from two independent registrations.} The first varied the
\emph{descriptor family}: ORB, AKAZE and mutual information all rank L1-only ahead of
adversarial+L1 on both the Ankara and European sets, with ORB at $-0.613 \pm 0.135$~px.
\textbf{Two caveats travel with these numbers and are stated rather than buried.} The ORB
figure is a paired difference over the \textbf{29 chips where both arms matched}, not over
the 53 chips where the L1-only arm matched at all; the AKAZE figure rests on 11 paired
chips. And the mutual-information margin of $-1.260 \pm 0.261$ is a \textbf{lower bound},
not an estimate: the registered subpixel refinement never ran, and the search grid censors
15.8\% of chips at its bound, censoring the worse arms harder and therefore compressing the
margin toward zero.

The second registration varied the \emph{matcher family} rather than the descriptor: KLT,
an NCC template grid, and phase correlation, crossed with two band conversions and an
urban subset over 6,510 scored comparisons. \textbf{The arm ordering is preserved in every
condition cell except two}, both at the European urban subset under phase correlation,
one in each band conversion, and both far below the threshold at which an ordering change
would be claimed. A registered prediction that the template matcher would favour the
sharper arm \textbf{failed in the direction that strengthens the result}: the L1-only
arm's margin grew rather than shrank, in all eight set-by-conversion combinations.

\textbf{The reading.} The candidate is refuted by two registrations that share no matcher,
no chip set and no scoring code. \textbf{That is stronger evidence than either alone}, and
it is the reason this row is reported as two tests rather than one.

\subsection{Corrected georeferencing in the fine-tuning pairs}

\textbf{The objection.} The fine-tuning pairs may be better georeferenced than the data
the published generator saw, so the advantage would come from the training data rather
than from the loss.

\textbf{The answer is in the design and needs no measurement.} No registered positional
contrast compares a fine-tuned arm with the pretrained generator: all four arms are
fine-tuned on identical pairs, so any georeferencing improvement is common to them and
cancels.

\textbf{Why this replaces a measurement we previously reported.} An earlier decomposition
attributed roughly 86\% of the European gain to scatter reduction rather than to a
systematic shift, and was reported as refuting this candidate. \textbf{That measurement's
per-chip artifacts did not survive and it cannot be reproduced}, which is disclosed in the
data-availability statement. The design argument is not a fallback: \textbf{it is
stronger, because it depends on the structure of the experiment rather than on a
computation, and nothing can be lost that would make it unverifiable.}

\subsection{Why the result is matcher-independent, mechanically}

A blurred template gives a broad correlation peak; invented structure gives a sharp peak
in the wrong place. \textbf{A broad peak in the right place localises better than a sharp
peak in the wrong one}, and that is true of any matcher that localises by correlation.
This is offered as the mechanical reason the previous row's result is unsurprising, not
as further evidence for it.
